%===========================================================
% Sección de Jiménez Reyes Alan Yael: bitácora clase 11 (IPC)
%===========================================================
\subsection{Clase 11: Comunicación entre procesos}
\textbf{Autor:} Jiménez Reyes Alan Yael

\textbf{Planeado:} RetroArch expone \texttt{/tmp/retroarch\_telemetry.sock} (\texttt{AF\_UNIX}, \texttt{SOCK\_DGRAM}) y un script de telemetría se conecta a él.

\textbf{Construido:}
\begin{itemize}
    \item \texttt{telemetria.py} pide la RAM del juego a RetroArch por UDP en \texttt{127.0.0.1:55355} (línea 50).
    \item Envía cada orden como datagrama a \texttt{/tmp/arcade\_act.sock} (línea 61).
    \item \texttt{puente\_hid.py} recibe el datagrama (socket en la línea 73), lo traduce a un reporte de salida HID de 8 bytes (línea 46) y lo escribe en \texttt{/dev/arcade\_hid} (línea 105)[cite: 23].
\end{itemize}

\textbf{Direcciones de RAM usadas:} Pac-Man (core Nestopia), dirección \texttt{0x067}, significa el contador de vidas del jugador (empieza en 3 y baja 1 con cada muerte).

\textbf{Diferencias:} RetroArch no ofrece un socket UNIX, sino un servidor de comandos por UDP en el puerto 55355; por lo tanto, el socket UNIX quedó implementado de manera interna como un puente entre nuestros dos procesos de Python. En el Avance 3 primero se usó un serial CDC para las órdenes; se cambió a reportes de salida HID para respetar el diagrama (un solo enlace USB HID)[cite: 23].

\subsection{Declaración de uso de IA}
Para la elaboración de este avance, utilicé asistencia de Inteligencia Artificial para verificar la correcta ejecución de los commits solicitados y estructurar los textos de esta bitácora.
